\section{Preliminaries}\label{sec:prelims}
Throughout, we use $\N$ to denote the set of natural numbers, and $\N_0$ for the set of non-negative integers. 
As stated in the introduction, for a finite group $G$ we use $\Irr(G)$ to denote a complete set of the ordinary irreducible characters of $G$. Further, we use $\Ch(G)$ to denote the set of all ordinary characters of $G$, and $\triv_G$ to mean the trivial character of $G$ (omitting the subscript when the meaning is clear from context).
For a subgroup $H\le G$ and $\phi\in\Irr(H)$, we let $\Irr(G\mid \phi)$ denote the set of $\chi\in\Irr(G)$ such that the restriction of $\chi$ to $H$ contains $\phi$ as a constituent. 

For $g\in G$ and $H\le G$, let $H^g\coloneqq gHg^{-1}$. Given $\phi\in\Ch(H)$, the character $\phi_g\in\Ch(H^g)$ is defined by $\phi_g(x)=\phi(g^{-1}xg)$. Mackey's Theorem allows us to describe restrictions and inductions between  subgroups of a finite group (see \cite[Chapter 5]{I}, for example).
\begin{theorem}[Mackey]\label{thm:Mackey}
	Let $G$ be a finite group and $H,K\le G$. Let $\phi\in\Ch(H)$. Then
	\[ \phi\up_H^G\down_K = \sum_{g\in K\setminus G/H} \phi_g \up_{H^g\cap K}^K, \]
	where the sum runs over a set of $(K,H)$-double coset representatives.
\end{theorem}



\subsection{Representation theory of symmetric groups}\label{sec:prelims-sn}
Next, we recall some key facts concerning the representation theory of symmetric groups, and refer the reader to \cite{J,JK} for further detail.
It is well known that $\Irr(S_n)$ is naturally in bijection with the set $\cP(n)$ of all partitions of $n$. By convention, $\cP(0)=\{\emptyset\}$ where $\emptyset$ denotes the empty partition. 
The irreducible character of $S_n$ corresponding to the partition $\lambda\vdash n$ will be denoted by $\chi^\lambda$. In particular, $\chi^{(n)}=\triv_{S_n}$, the trivial character of $S_n$, and $\chi^{(1^n)}=\sgn_{S_n}$, the sign character of $S_n$. If $\alpha$ is a (finite) sequence of integers but is not a partition, then we interpret $\chi^\alpha$ to be the zero function.

The Young diagram of a partition $\lambda$ will be denoted by $[\lambda]$, and that of a skew shape $\lambda/\mu$ by $[\lambda/\mu]\coloneqq [\lambda]\setminus[\mu]$ for $\mu$ a subpartition of $\lambda$ (written $\mu\subseteq\lambda$). The boxes in a Young diagram will sometimes be referred to as nodes, and we refer to skew shapes and skew diagrams interchangeably when the meaning is clear.
We denote the length of the partition $\lambda$ by $l(\lambda)$, and the conjugate partition of $\lambda$ by $\lambda'$. Note
\begin{equation}\label{eqn:sign}
	\chi^{\lambda'}=\sgn_{S_n}\cdot\ \chi^\lambda
\end{equation}
for all $\lambda\vdash n$ (see \cite[2.1.8]{JK}). 

We record some operations on partitions. Let $\lambda=(\lambda_1,\lambda_2,\dotsc)$ and $\mu=(\mu_1,\mu_2,\dotsc)$ be two partitions. Then $+$ denotes component-wise addition, i.e.~$\lambda+\mu=(\lambda_1+\mu_1,\lambda_2+\mu_2,\dotsc)$, and $\lambda\sqcup\mu$ denotes the partition obtained by taking the disjoint union of the parts of $\lambda$ and $\mu$ and reordering so that parts are in non-increasing order.
When clear from context, we abbreviate $(a^b)\coloneqq (a,a,\dotsc,a)$ where there are $b$ parts of size $a$; in general we will specify $(a^b)\vdash ab$ to avoid confusion with the single part partition of $a^b$.

We also make use of skew characters for symmetric groups, i.e.~those indexed by skew shapes. For partitions $\mu$ and $\lambda$ such that $|\mu|\le|\lambda|$, the skew character $\chi^{\lambda/\mu}$ of $S_{|\lambda|-|\mu|}$ satisfies
\begin{equation}\label{eqn:skew}
	\langle \chi^{\lambda/\mu},\chi^\nu\rangle = \langle \chi^\lambda, (\chi^\mu\times \chi^\nu)\up_{S_{|\mu|}\times S_{|\nu|}}^{S_{|\mu|+|\nu|}}\rangle\qquad\forall\ \nu\vdash |\lambda|-|\mu|.
\end{equation}
Note if $\mu\not\subseteq\lambda$ then $\chi^{\lambda/\mu}=0$. This can be seen from the Littlewood--Richardson rule, which gives an explicit combinatorial description of the decomposition into irreducibles of the induced character appearing in the above expression, with Littlewood--Richardson coefficients arising as the multiplicities. These appear in many contexts, so we shall now fix the notation which will be used throughout this article (see \cite{J}).

\begin{definition}
	Let $n\in\N$, $\lambda=(\lambda_1,\dotsc,\lambda_k)\vdash n$ and $\mathcal{C}=(c_1,c_2,\dotsc,c_n)$ be a sequence of positive integers. We say that $\mathcal{C}$ is of type $\lambda$ if 
	\[ |\{i\in\{1,2,\dotsc,n\} \mid c_i=j\}|=\lambda_j\]
	for all $j\in\{1,2,\dotsc,k\}$. Moreover, we say that an element $c_j$ of $\mathcal{C}$ is good if $c_j=1$ or if
	\[ |\{i\in\{1,2,\dotsc,j-1\}\mid c_i=c_j-1\}| > |\{i\in\{1,2,\dotsc,j-1\}\mid c_i=c_j\}|. \]
	Finally, we say that $\mathcal{C}$ is good if $c_j$ is good for all $j\in\{1,\dotsc,n\}$.
\end{definition}

\begin{theorem}[Littlewood--Richardson rule]\label{thm:LR}
	Let $m,n\in\mathbb{N}_0$. Let $\mu\vdash m$ and $\nu\vdash n$. Then
	\[ (\chi^\mu\times\chi^\nu)\up^{S_{m+n}}_{S_m\times S_n} = \sum_{\lambda\vdash m+n} c^\lambda_{\mu,\nu}\cdot \chi^\lambda\]
	where the Littlewood--Richardson coefficient $c^\lambda_{\mu,\nu}$ equals the number of ways to replace the nodes of the skew Young diagram of $\lambda/\mu$ by natural numbers such that
	\begin{enumerate}[label=(\roman*)]
		\item\label{it1} the sequence obtained by reading the numbers from right to left, top to bottom is good of type $\nu$;
		\item \label{it2} the numbers are non-decreasing \textup{(}weakly increasing\textup{)} left to right along rows; and
		\item \label{it3} the numbers are strictly increasing down columns.
	\end{enumerate}
\end{theorem}

We call the order in Theorem~\ref{thm:LR}\ref{it1} the reading order of a skew shape. Let $\nu$ be a partition and $\gamma$ be a skew shape of size $|\nu|$. We call a way of replacing the nodes of $\gamma$ by numbers satisfying conditions Theorem~\ref{thm:LR}\ref{it1}--\ref{it3} a Littlewood--Richardson (LR) filling of $\gamma$ of type $\nu$. 
Clearly $c^\lambda_{\mu,\nu}=c^\lambda_{\nu,\mu}$. Using Littlewood--Richardson coefficients, we can also rephrase \eqref{eqn:skew} as $\chi^{\lambda/\mu}=\sum_\nu c^\lambda_{\mu,\nu}\cdot \chi^\nu$. Moreover, we can extend this notation to generalised Littlewood--Richardson coefficients $c^\lambda_{\mu^1,\dotsc,\mu^r}$ describing the constituents of
\[ (\chi^{\mu^1}\times\cdots\times\chi^{\mu^r})\up_{S_{n_1}\times\cdots\times S_{n_r}}^{S_{n_1+\cdots+n_r}}, \]
for any $r\in\N$ and $n_i\in\N_0$, and partitions $\mu^i\vdash n_i$ and $\lambda\vdash n_1+\cdots+n_r$. Similarly, $c^\lambda_{\mu^1,\dotsc,\mu^r}=c^\lambda_{\mu^{\sigma(1)},\dotsc,\mu^{\sigma(r)}}$ for any $\sigma\in S_r$. Furthermore, for $A\subseteq\cP(n)$ and $B\subseteq\cP(m)$ we define the operation $\star$ as follows:
\[ A\star B\coloneqq  \{\lambda\vdash m+n \mid \exists\ \mu\in A,\ \nu\in B\text{ s.t. }c^\lambda_{\mu,\nu}>0 \}. \]
We note that $\star$ is both commutative and associative.

Finally, for a partition $\lambda=(\lambda_1,\lambda_2,\dotsc,\lambda_k)\vdash n$, we let $S_\lambda\cong S_{\lambda_1}\times\cdots\times S_{\lambda_k}$ denote the corresponding Young subgroup of $S_n$. The permutation module $\triv_{S_\lambda}\up^{S_n}$ induced by the action of $S_n$ on the cosets of $S_\lambda$ will be denoted by $M^\lambda$. 
Young's Rule (see \cite[2.8.5]{JK}) tells us the decomposition of these permutation modules into irreducibles. Denoting the character of $M^\lambda$ by $\xi^\lambda$, we have that $\langle \xi^\lambda,\chi^\alpha\rangle$ equals the number of semistandard Young tableaux of shape $\alpha$ and content $\lambda$, for any $\alpha\vdash n$. Moreover, this multiplicity is positive if and only if $\alpha\trianglerighteq\lambda$, where $\trianglerighteq$ denotes the dominance partial order on partitions.



\subsection{Wreath products and Sylow subgroups of symmetric groups}
In order to describe the Sylow subgroups of symmetric groups, we briefly introduce some notation for wreath products. Let $G$ be a finite group and let $H\le S_n$ for some $n\in\N$. The natural action of $S_n$ on the factors of the direct product $G^{\times n}$ induces an action of $S_n$ (and therefore of $H$) via automorphisms of $G^{\times n}$, giving the wreath product $G\wr H\coloneqq  G^{\times n}\rtimes H$.  
As in \cite[Chapter 4]{JK}, we denote the elements of $G\wr H$ by $(g_1,\dotsc,g_n;h)$ for $g_i\in G$ and $h\in H$. Let $V$ be a $\bC G$--module and suppose it affords the character $\phi$. Let $V^{\otimes n}$ be the corresponding $\bC G^{\times n}$--module. The left action of $G\wr H$ on $V^{\otimes n}$ defined by linearly extending
\[ (g_1,\dotsc,g_n;h)\ :\quad v_1\otimes \cdots\otimes v_n \longmapsto g_1v_{h^{-1}(1)}\otimes\cdots\otimes g_nv_{h^{-1}(n)} \]
turns $V^{\otimes n}$ into a $\bC(G\wr H)$--module, which we denote by $\widetilde{V^{\otimes n}}$ (see \cite[(4.3.7)]{JK}), and we denote its character by $\tilde{\phi}$. For any $\psi\in\Ch(H)$, we define $\cX(\phi;\psi)$ as follows:
\[ \cX(\phi;\psi)\coloneqq \tilde{\phi}\cdot\Infl_H^{G\wr H}(\psi)\ \in\Ch(G\wr H). \]
The inflation $\Infl_H^{G\wr H}(\psi)$ of $\psi$ from $H$ to $G\wr H$ (identifying $H$ with the quotient $(G\wr H)/G^{\times n}$) is sometimes abbreviated to simply $\psi$, for convenience.

\begin{lemma}\label{lem:infl-ind}
	Let $G$, $H$ and $\phi$ be as above. Let $L\le H$ and $\tau\in\Ch(L)$. Then $\cX(\phi;\tau)\up_{G\wr L}^{G\wr H} = \cX(\phi;\tau\up_L^H)$.
\end{lemma}
\begin{proof}
	For any $\alpha\in\Ch(H)$ and $\beta\in\Ch(L)$, it is easy to check that $\alpha\cdot(\beta\up^H)=(\alpha\down_L\cdot \beta)\up^H$. Hence
	\[ \cX(\phi;\tau)\up_{G\wr L}^{G\wr H} \coloneqq  \big( \tilde{\phi}\down_{G\wr L}^{G\wr H}\cdot\Infl_L^{G\wr L}\tau\big)\up_{G\wr L}^{G\wr H} = \tilde{\phi} \cdot \big(\Infl_L^{G\wr L}(\tau)\big)\up_{G\wr L}^{G\wr H}. \]
	But induction and inflation of characters commute, so $\big(\Infl_L^{G\wr L}(\tau)\big)\up_{G\wr L}^{G\wr H} = \Infl_H^{G\wr H}(\tau\up_L^H)$. Thus $\cX(\phi;\tau)\up_{G\wr L}^{G\wr H} = \tilde{\phi}\cdot \Infl_H^{G\wr H}(\tau\up_L^H) = \cX(\phi;\tau\up_L^H)$, as claimed.
\end{proof}

Furthermore, if $\phi\in\Irr(G)$ then Gallagher's Theorem \cite[Corollary 6.17]{I} gives $\Irr(G\wr H\mid \phi^{\times n}) = \{\cX(\phi;\psi)\mid \psi\in\Irr(H)\}$, where $\Irr(G\wr H\mid \phi^{\times n})\coloneqq \{\chi\in\Irr(G\wr H)\mid \langle \chi\down_{G^{\times n}}, \phi^{\times n}\rangle \ne 0 \}$.
For a full description of $\Irr(G\wr H)$, we refer the reader to \cite[Chapter 4]{JK}; in the case $H=S_2$, we use the notation below.

\begin{notation}\label{not:GwrS2}
	Let $G$ be a finite group and suppose $\Irr(G)=\{\chi^i \mid i\in I\}$. Then 
	\[ \Irr(G\wr S_2) = \{ \psi^{i,j} \mid i\ne j\in I \} \sqcup \{ \psi^{i,i}_+, \psi^{i,i}_- \mid i\in I \} \]
	where 
	\[ \psi^{i,j} \coloneqq  (\chi^i\times\chi^j)\up_{G\times G}^{G\wr S_2}\ = \psi^{j,i},\quad \psi^{i,i}_{+} \coloneqq  \cX(\chi^i;\triv_{S_2})\quad\text{and}\quad \psi^{i,i}_{-} \coloneqq  \cX(\chi^i;\sgn_{S_2}). \]
\end{notation}

It will be useful to describe the decomposition of the permutation character $\triv_{H\wr S_2}\up^{G\wr S_2}$, for finite groups $H\le G$.

\begin{lemma}\label{lem:Prop1.2}
	Let $G$ be a finite group and let $\Irr(G)$ and $\Irr(G\wr S_2)$ be as in Notation~\ref{not:GwrS2}.
	Let $H\leq G$ and let $\pi \coloneqq  \triv_{H}\up^{G}$. Then $\tilde{\pi} \coloneqq  \triv_{H\wr S_2}\up^{G\wr S_2}$ decomposes into irreducible constituents with multiplicities given by
	\[ \langle\tilde{\pi},\psi^{i,j}\rangle = \langle\pi,\chi^i\rangle \cdot \langle\pi,\chi^j\rangle \quad\text{and}\quad \langle\tilde{\pi},\psi^{i,i}_{\pm}\rangle = \tfrac{1}{2}\cdot\big( \langle\pi,\chi^i\rangle^2 \pm \langle\pi,\chi^i\rangle \big). \]
\end{lemma}

\begin{proof}
	The first part follows from Mackey's theorem applied to the subgroups $G\times G$ and $H\wr S_2$ of $G\wr S_2$: since $(G\times G)\cdot (H\wr S_2)=G\wr S_2$ and $(G\times G)\cap (H\wr S_2)=H\times H$, we have
	\[ \langle\tilde{\pi},\psi^{i,j}\rangle	= \langle \triv_{H\wr S_2}\up^{G\wr S_2}\down_{G\times G}, \chi^i\times\chi^j\rangle = \langle \triv_{H\times H}\up^{G\times G}, \chi^i\times \chi^j\rangle = \langle\pi,\chi^i\rangle \cdot \langle\pi,\chi^j\rangle. \]
	Next, we use the wreath product character formula \cite[Lemma 4.3.9]{JK} to obtain
	\[ \psi^{i,i}_{\pm}(g_1,g_2;1) = \chi^i(g_1)\cdot\chi^i(g_2)\quad\text{and}\quad\psi^{i,i}_{\pm}\big(g_1,g_2;(1,2)\big) = \pm\chi^i(g_1 g_2) \qquad\forall\ g_1, g_2\in G. \]
	Hence
	\begin{align*}
		\langle\tilde{\pi},\psi^{i,i}_{\pm}\rangle &= \langle\triv_{H\wr S_2},\psi^{i,i}_{\pm}\down_{H\wr S_2}\rangle\\
		&= \frac{1}{|H\wr S_2|}\sum_{h_1,h_2\in H}\big(\chi^i(h_1)\cdot\chi^i(h_2)\pm\chi^i(h_1 h_2)\big) = \tfrac{1}{2}\big(\langle\pi,\chi^i\rangle^2\pm \langle\pi,\chi^i\rangle\big)
	\end{align*}
	as claimed.
\end{proof}

To describe Sylow subgroups of symmetric groups, fix a prime $p$ and let $n\in\N$. Let $P_n$ denote a Sylow $p$-subgroup of $S_n$. Clearly $P_1$ is trivial while $P_p$ is cyclic of order $p$. More generally, $P_{p^k}= (P_{p^{k-1}})^{\times p}\rtimes P_p=P_{p^{k-1}}\wr P_p\cong P_p\wr \cdots \wr P_p$ ($k$-fold wreath product) for all $k\in\N$. Now suppose $n\in\N$ and let $n=\sum_{i=1}^k a_ip^{n_i}$ be its $p$-adic expansion, i.e.~$n_1>\cdots>n_k\ge 0$ and $a_i\in\{1,2,\dotsc,p-1\}$ for all $i$. Then $P_n\cong (P_{p^{n_1}})^{\times a_1}\times\cdots\times (P_{p^{n_k}})^{\times a_k}$. 

To fix a convention for denoting such wreath products involving Sylow subgroups of symmetric groups more generally, we have the following.

\begin{notation}\label{not:convention}
	Let $G$ be a finite group and $p$ a prime. We use the convention that $P_n$ will always be viewed as a subgroup of $S_n$ in the notation $G\wr P_n$, that is, $G\wr P_n$ is a semi-direct product $G^{\times n}\rtimes P_n$. 
\end{notation}

\begin{remark}\label{rem:convention}
	Suppose $p=2$ and $n=2^{n_1}+\cdots+2^{n_k}$ for some $n_1>\cdots>n_k\ge 0$. With the convention of Notation~\ref{not:convention}, we observe that 
	\begin{align*}
		P_{2n} \cong P_{2^{n_1+1}}\times\cdots\times P_{2^{n_k+1}} &\cong (P_2\wr P_{2^{n_1}})\times\cdots\times (P_2\wr P_{2^{n_k}})\\
		&\cong P_2 \wr (P_{2^{n_1}}\times \cdots \times P_{2^{n_k}})\cong P_2 \wr P_n,
	\end{align*}
	viewing $P_{2^{n_1}}\times \cdots \times P_{2^{n_k}}\cong P_n$ naturally as a subgroup of $S_n$. Inductively, we also have $P_{2^tn}\cong P_{2^t}\wr P_n$ for all $t\in\N$.
	On the other hand, we clarify for example that $P_2\wr P_3 \not\cong P_2\wr P_2$ in this notation, even though $P_3\cong P_2$.%\hfill$\lozenge$
\end{remark}

We now return to an arbitrary prime $p$. Following the notation introduced in \cite{GL2}, given $\chi\in\Irr(S_n)$ and $\phi\in\Irr(P_n)$, the \textit{Sylow branching coefficient} $Z^\chi_\phi$ denotes the non-negative integer
\[ Z^\chi_\phi\coloneqq \langle \chi\down_{P_n}, \phi\rangle. \]
In this article, we will be particularly interested in the case where $\phi=\triv_{P_n}$, and abbreviate $Z^\chi_{\triv_{P_n}}$ to $Z^\chi$. Moreover, if $\chi=\chi^\lambda$ for a partition $\lambda$, then we shorten $Z^{\chi^\lambda}$ to $Z^\lambda$. 

We record one more lemma which will be useful later.
\begin{lemma}\label{lem:X-induced}
	Let $A$ and $B$ be finite groups, and let $n\in\N$. Then
	\[ \triv\up_{(A\times B)\wr S_n}^{A\wr S_n \times B\wr S_n} = \sum_{\phi\in\Irr(S_n)} \cX(\triv_A;\phi)\cdot\cX(\triv_B;\phi). \]
\end{lemma}
\begin{proof}
	Let $\phi\in\Irr(S_n)$. By Frobenius reciprocity, 
	\begin{align*}
		\langle \triv\up_{(A\times B)\wr S_n}^{A\wr S_n \times B\wr S_n}, \cX(\triv_A;\phi)\cdot\cX(\triv_B;\phi)\rangle &= \langle \triv, \big( \cX(\triv_A;\phi)\cdot\cX(\triv_B;\phi) \big) \down_{(A\times B)\wr S_n}\rangle\\
		&= \tfrac{1}{|(A\times B)\wr S_n|}\sum \cX(\triv_A;\phi)(x)\cdot\cX(\triv_B;\phi)(x),
	\end{align*}
where the sum runs over all $x\in (A\times B)\wr S_n$.
	But this equals $\tfrac{|A|^n\cdot|B|^n}{|(A\times B)\wr S_n|}\sum_{g\in S_n} \phi(g)^2$ by \cite[Lemma 4.3.9]{JK}. As symmetric group characters are real-valued (in fact, integer-valued), this simplifies to $\tfrac{1}{|S_n|}\sum_{g\in S_n}\overline{\phi(g)}\cdot\phi(g)=\langle \phi,\phi\rangle=1$.
\end{proof}



\subsection{Plethysms and deflations}\label{sec:plethysm-deflation}
When $\phi$ and $\psi$ are characters of symmetric groups, the characters $\cX(\phi;\psi)$ introduced above are closely related to plethysms of Schur functions: we give a brief description here and refer the reader to \cite{dBPW, Mac, Stanley} for further detail. Let $s_\lambda$ denote the Schur function corresponding to the partition $\lambda$, and $\circ$ the plethystic product of symmetric functions. Using the characteristic map (see e.g.~\cite[Chapter 7]{Stanley}) between class functions of finite symmetric groups and the ring of symmetric functions, we have the correspondence
\[ \cX(\chi^\nu;\chi^\lambda)\up_{S_{|\nu|}\wr S_{|\lambda|}}^{S_{|\nu|\cdot|\lambda|}}\ \longleftrightarrow\ s_\lambda\circ s_\nu \]
for all partitions $\lambda$ and $\nu$. Therefore the plethysm coefficient $a^\mu_{\lambda,\nu}$ satisfies
\begin{equation}\label{eqn:a}
	a^\mu_{\lambda,\nu} = \langle s_\lambda\circ s_\nu, s_\mu\rangle = \langle \cX(\chi^\nu;\chi^\lambda)\up_{S_{|\nu|}\wr S_{|\lambda|}}^{S_{|\nu|\cdot|\lambda|}}, \chi^\mu \rangle	
\end{equation}
for all partitions $\mu$ (and note that this is zero if $|\mu|\ne|\nu|\cdot|\lambda|$). 
We also introduce plethysm coefficients indexed by skew shapes: for partitions $\beta\subseteq\alpha$, $\delta\subseteq\gamma$ and $\nu$,
\begin{equation}\label{eqn:skew-plethysm}
	a^{\alpha/\beta}_{\gamma/\delta,\nu}\coloneqq  \sum_{\substack{\eta\vdash|\gamma|-|\delta|\\\zeta\vdash|\alpha|-|\beta|}} c^\gamma_{\eta,\delta}\cdot c^\alpha_{\zeta,\beta}\cdot a^\zeta_{\eta,\nu}.
\end{equation}
In other words, if $\phi$ and $\theta$ are skew shapes and $\nu$ is any partition, then $a^\phi_{\theta,\nu} = \langle \cX(\chi^\nu; \chi^\theta)\up_{S_{|\nu|}\wr S_{|\theta|}}^{S_{|\nu|\cdot|\theta|}}, \chi^\phi\rangle$, extending the equality in \eqref{eqn:a}. We also remark that if $\zeta$ and $\eta$ are partitions, then $a^\zeta_{\eta,\emptyset}=1$ if $\zeta=\emptyset$ and $\eta=(n)$ for some $n$, and $a^\zeta_{\eta,\emptyset}=0$ otherwise. 

A well known symmetry property of plethysm coefficients involving the conjugation involution is the following (see e.g.~\cite[Ex.~1, Ch. I.8]{Mac}):
\begin{lemma}\label{lem:conj-plethysm}
	Let $\lambda$, $\mu$ and $\nu$ be partitions. Then
	\[ a^\nu_{\lambda,\mu} = a^{\nu'}_{\lambda^\ast,\mu'},\quad\text{where}\quad \lambda^\ast\coloneqq \begin{cases} \lambda & \text{if $|\mu|$ is even},\\ \lambda' & \text{if $|\mu|$ is odd}.\end{cases} \]
\end{lemma}

Character deflations were introduced in \cite[Definition 1.1]{EPW}, and used to prove results generalising the Murnaghan--Nakayama rule for computing symmetric group character values and (special cases of) the Littlewood--Richardson rule, as well as to verify new cases of the long-standing Foulkes' Conjecture. We observe that they give another language in which to describe certain plethysm coefficients. We first record the definition of these deflations in notation which we have introduced thus far.
\begin{definition}\label{def:deflation}
	Let $m,n\in\N$ and $\theta\in\Irr(S_m)$. Let $\xi\in\Irr(S_m\wr S_n)$. Then 
	\[ \operatorname{Def}^\theta_{S_n}(\xi) \coloneqq  \begin{cases}
		\chi^\nu & \text{ if }\xi = \cX(\theta;\chi^\nu)\text{ for some }\nu\vdash n,\\
		0 & \text{ otherwise},
	\end{cases} \]
	which then extends linearly to all class functions of $S_m\wr S_n$. If $\chi$ is a class function of $S_{mn}$ then set
	\[ \operatorname{Defres}^\theta_{S_n}(\chi) \coloneqq  \operatorname{Def}^\theta_{S_n}(\chi\down_{S_m\wr S_n}). \]
	When $n\in\N$ is fixed and $\gamma\vdash mn$ for some $m\in\N$, we use the notation
	\[ \delta^\gamma\coloneqq \operatorname{Defres}^{\triv_{S_m}}_{S_n}(\chi^\gamma) = \operatorname{Def}^{\triv_{S_m}}_{S_n}(\chi^\gamma\down^{S_{mn}}_{S_m\wr S_n}). \]
	In this article, we will sometimes refer to $\delta^\gamma$ as the deflation of $\gamma$ with respect to $S_n$ \textup{(}where the $m$ is understood from $|\gamma|/n$ and we suppress $m$ from the notation\textup{)}.
\end{definition}
In particular, if $\lambda\vdash n$ then we have that $a^\gamma_{\lambda,(m)} = \langle \delta^\gamma,\chi^\lambda \rangle$. This relation between plethysm coefficients and deflations can immediately be extended to skew shapes using Littlewood--Richardson coefficients. Namely, if $\alpha,\beta,\nu,\lambda$ are partitions such that $|\alpha|-|\beta|=n$ and $|\nu|-|\lambda|=mn$, then we may define $\delta^{\nu/\lambda}\coloneqq \operatorname{Defres}^{\triv_{S_m}}_{S_n}(\chi^{\nu/\lambda}) = \operatorname{Def}^{\triv_{S_m}}_{S_n}(\chi^{\nu/\lambda}\down^{S_{mn}}_{S_m\wr S_n})$. Then $\chi^{\alpha/\beta}=\sum_\sigma c^\alpha_{\beta,\sigma}\cdot \chi^\sigma$ and $\delta^{\nu/\lambda}=\sum_\tau c^\nu_{\lambda,\tau}\cdot \delta^\tau$ give
\begin{equation}\label{eqn:plethysm-deflation-skew}
	a^{\nu/\lambda}_{\alpha/\beta,(m)} = \langle \delta^{\nu/\lambda}, \chi^{\alpha/\beta}\rangle.
\end{equation}

The following is a straightforward result on plethysm coefficients involving `tall' partitions. We include a proof for convenience.
\begin{lemma}\label{lem:tall-pleth}
	Let $m,n\in\N$, $\lambda\vdash mn$ and $\nu\vdash n$. If $l(\lambda)>n$, then $a^\lambda_{\nu,(m)}=0$.
\end{lemma}

\begin{proof}
	Let $a\coloneqq a^\lambda_{\nu,(m)}= \langle \chi^\lambda\down^{S_{mn}}_{S_m\wr S_n}, \cX(\chi^{(m)};\chi^\nu)\rangle$. Let $Y\coloneqq (S_m)^{\times n}\le S_m\wr S_n$. 
	Let $\cX\coloneqq \cX(\chi^{(m)};\chi^\nu)$ and write $\chi^\lambda\down_{S_m\wr S_n} = a\cX + \Delta$ for some $\Delta\in\Ch(S_m\wr S_n)$. Since $\cX\down_Y = \chi^\nu(1)\cdot(\chi^{(m)})^n=\chi^\nu(1)\cdot\triv_Y$, it follows that $\langle \chi^\lambda\down_Y,\triv_Y\rangle \ge \langle a\cdot \cX\down_Y, \triv_Y\rangle = a\cdot\chi^\nu(1).$
	But by the Littlewood--Richardson rule, $l(\lambda)>n$ implies $\langle \chi^\lambda\down_Y,\triv_Y\rangle=0$ since $\triv_Y = (\chi^{(m)})^n$, which gives $a=0$.
\end{proof}

Finally, we record the following result of Thrall \cite{T}. We note that a partition is even if all of its parts are of even size.
\begin{proposition}\label{prop:thrall}
	Let $n\in\N$. Then 
	\begin{enumerate}[label=\textup{(\roman*)}]
		\item $s_{(n)}\circ s_{(2)} = \sum_\lambda s_\lambda$ where the sum runs over all even partitions $\lambda\vdash 2n$, and
		\item $s_{(2)}\circ s_{(n)} = \sum_\mu s_\mu$ where the sum runs over all even partitions $\mu\vdash 2n$ with at most two parts.
	\end{enumerate}
\end{proposition}
In other words, $a^\lambda_{(n),(2)}=1$ if $\lambda\vdash 2n$ is even and $a^\lambda_{(n),(2)}=0$ otherwise, while $a^\mu_{(2),(n)}=1$ if $\mu\vdash 2n$ is even and $l(\mu)\le 2$, and $a^\mu_{(2),(n)}=0$ otherwise.



